\documentclass[a4paper,11pt]{article}

\usepackage[utf8]{inputenc}
\usepackage[T2A]{fontenc}
\usepackage[russian]{babel}
\usepackage{listings}
\usepackage{geometry}
\usepackage{hyperref}
\geometry{margin=2.5cm}

\lstset{
  basicstyle=\ttfamily\small,
  showstringspaces=false,
  breaklines=true,
  frame=single,
  framesep=4pt
}

\title{rjson: расширение JSON метками и ссылками}
\author{Кабиров Рустам, Кабиров Руслан, Заходов Егор\\ МФТИ, 2026}
\date{}

\begin{document}
\maketitle

\section{Постановка}

Формат rjson расширяет JSON двумя конструкциями:
\begin{itemize}
  \item метка \texttt{ident@value} — даёт значению имя;
  \item ссылка \texttt{@ident} — в позиции значения раскрывается в
        значение соответствующей метки.
\end{itemize}
Дополнительно разрешены неэкранированные ключи объектов
(\texttt{name:} вместо \texttt{"name":}). Компилятор разрешает все
ссылки и выводит чистый JSON.

\section{Грамматика (PEG)}

Используем формализм PEG: \texttt{/} — упорядоченный выбор,
\texttt{?},\texttt{*},\texttt{+} — кванторы.
\begin{lstlisting}
Value      <- _ (IdentStart / Ref / Object / Array / String / Number) _
IdentStart <- Ident ("@" Value / Keyword)
Ref        <- "@" Ident
Object     <- "{" (Pair ("," Pair)*)? "}"
Pair       <- (String / Ident) ":" Value
Array      <- "[" (Value ("," Value)*)? "]"
Ident      <- [a-zA-Z_] [a-zA-Z_0-9]*
Keyword    <- "true" / "false" / "null"
String     <- '"' (Escape / [^"\])* '"'
Number     <- "-"? [0-9]+ ("." [0-9]+)? ([eE] [+-]? [0-9]+)?
_          <- (whitespace)*
\end{lstlisting}

\section{Парсер}

Реализован на библиотеке Parsec в монадическо-аппликативном стиле, де можно — аппликативно, где нужна зависимость — через do.

Каждая лексема съедает хвостовые пробелы.

Сложность представляют идентификаторы: они начинают сразу три
конструкции — метку (\texttt{ident@\ldots}), ключ объекта
(\texttt{ident:\ldots}) и ключевое слово (\texttt{true}, \texttt{false},
\texttt{null}). Различение по контексту:

\begin{itemize}
  \item внутри \texttt{Value} мы парсим идентификатор и затем смотрим
        вперёд: если идёт \texttt{@}, это метка; иначе — ключевое
        слово (если совпадает) или ошибка;
  \item внутри \texttt{Pair} за идентификатором ожидается \texttt{:}.
\end{itemize}

Такой алгоритм даёт однозначный детерминированный разбор без
комбинатора \texttt{try}.

\section{Разрешение ссылок}

Двухпроходный алгоритм.

\paragraph{Проход 1. Сбор таблицы меток.}
Обход дерева в глубину, накапливаем \texttt{Map String Value}: для
каждого узла \texttt{VLabel n inner} проверяем, не определена ли уже
метка $n$. Если определена — ошибка \texttt{Dup}. Иначе кладём $n
\mapsto inner$ и рекурсивно собираем метки из \texttt{inner}.

Поскольку таблица собирается \emph{до} подстановки, ссылка может
встречаться раньше определения метки (\emph{forward reference}).

\paragraph{Проход 2. Подстановка.}
Обход дерева с состоянием — список \texttt{path} активно
разворачиваемых имён:
\begin{itemize}
  \item \texttt{VLabel \_ inner} — снимаем обёртку;
  \item \texttt{VRef n} — если $n \in \texttt{path}$, выдаём \texttt{Cycle}
        c полным хронологическим путём; иначе находим $n$ в таблице
        (либо \texttt{Undef}) и продолжаем разрешение
        с $\texttt{path}' = n : \texttt{path}$;
  \item контейнеры — рекурсивно с тем же путём (это ловит вложенные
        циклы типа \texttt{a@\{x:@a\}}).
\end{itemize}

Псевдокод:
\begin{lstlisting}[language=Haskell]
subst table = go []
  where
    go path val = case val of
      VLabel _ inner -> go path inner
      VRef n
        | n `elem` path -> Left (Cycle (reverse path ++ [n]))
        | otherwise -> case Map.lookup n table of
            Nothing -> Left (Undef n)
            Just v' -> go (n : path) v'
      VObj ps -> VObj <$> traverse (\(k,x) -> ((,) k) <$> go path x) ps
      VArr xs -> VArr <$> traverse (go path) xs
      _       -> Right val
\end{lstlisting}

\section{Сериализация}

Отдельный ручной эмиттер выводит чистый JSON с двухпробельными
отступами; строки и числа форматируются стандартным \texttt{show}.

\section{Сборка и запуск}

\begin{lstlisting}[language=sh]
cabal build
cabal run rjson < examples/universities.rjson > out.json
\end{lstlisting}

Коды возврата: 0 — успех, 1 — ошибка парсинга, 2 — ошибка разрешения.


\end{document}
